%!TEX root = ../template.tex
%%%%%%%%%%%%%%%%%%%%%%%%%%%%%%%%%%%%%%%%%%%%%%%%%%%%%%%%%%%%%%%%%%%%
%% chapter3-architectures-draft.tex
%%
%% DRAFT: victim architectures under test. Intended to precede
%% Section~\ref{subsec:pgd_maddpg}.
%%%%%%%%%%%%%%%%%%%%%%%%%%%%%%%%%%%%%%%%%%%%%%%%%%%%%%%%%%%%%%%%%%%%

% ─────────────────────────────────────────────────────────────────────────────
\section{Victim Architectures Under Test}
\label{sec:victim_architectures}

The attacks are evaluated against the family of trained variants inherited from the
base system~\cite{amaral2021maddpg-routing} rather than against a single model, so
that conclusions about the routing task can be separated from conclusions about one
particular network design. All variants share the same actor: a multilayer
perceptron mapping the local observation to $n_d$ blocks of $K$ path scores, with
the executed route given by the $\arg\max$ within each block. They differ along
three binary axes, described below.

% ─────────────────────────────────────────────────────────────────────────────
\subsection{Central (CC) versus local (LC) critic}
\label{subsec:axis_critic_domain}

The two critic domains differ in what the value function is allowed to observe
during training:
\begin{equation}
  \underbrace{Q_{\phi^i}\big(\mathbf{s}, a_1, \dots, a_N\big)}_{\text{CC:\ } N(d_{\text{obs}} + |\mathcal{A}|)}
  \qquad\text{versus}\qquad
  \underbrace{Q_{\phi^i}\big(s^{\text{local}}_i, a_i\big)}_{\text{LC:\ } d_{\text{obs}} + |\mathcal{A}|},
  \label{eq:arch_critic_domains}
\end{equation}

\noindent with $d_{\text{obs}}$ the per-agent observation dimension and
$|\mathcal{A}| = n_d K$ the per-agent action dimension. The central critic is the
canonical MADDPG construction~\cite{lowe2017multi}: conditioning on the joint action
removes the non-stationarity each agent would otherwise face as its peers' policies
change, at the cost of an input that grows linearly in the number of agents. The
local critic reduces the scheme to independent actor--critic learners; it forfeits
that guarantee but is far cheaper and its size is independent of the agent count.

The distinction also determines which attack objectives are available. Because the
central critic is an explicit function of \emph{all} agents' actions, it supplies a
non-separable objective that couples the compromised agents, used in
Equation~\eqref{eq:obj_critic}. Local-critic variants admit no such objective and
are attacked with the congestion objective of Equation~\eqref{eq:obj_congestion}
instead.

% ─────────────────────────────────────────────────────────────────────────────
\subsection{Simple versus duelling value head}
\label{subsec:axis_value_head}

The simple head passes the concatenated state and action through two rectified
hidden layers and emits a scalar. The duelling head replaces the output layer with
two parallel streams whose outputs are summed, $Q(s,a) = V(h) + A(h)$.

One caveat should be stated. In the original duelling formulation the advantage
stream emits one value \emph{per action}, and subtracting the mean advantage makes
the split identifiable. Here the action is an \emph{input} to the critic rather than
an output index, so both streams emit scalars and no such constraint is applied: any
constant may be shifted between $V$ and $A$ without changing $Q$. The duelling head
is thus an extra parallel projection rather than a value--advantage separation, and
differences it produces should be attributed on that basis.
% TODO: cite Wang et al. (2016), Duelling Network Architectures, once added to the .bib

% ─────────────────────────────────────────────────────────────────────────────
\subsection{With and without GNN encoding}
\label{subsec:axis_gnn}

The GNN variants insert a shared encoder between the environment and the actor.
Each agent is a graph node whose features are its full observation, edges come from
the trainable subgraph of the physical topology, and two \texttt{GraphConv} layers
aggregate neighbour information under a residual connection:
\begin{equation}
  \tilde{\mathbf{X}} = \mathrm{ReLU}\Big(
     \mathrm{GraphConv}_2\big(\mathrm{ReLU}(\mathrm{GraphConv}_1(\mathbf{X}, E)), E\big)
     + \mathbf{X} \Big).
  \label{eq:arch_gnn}
\end{equation}

\noindent The encoder is dimension-preserving, so it is a drop-in pre-processor and
the actor is unchanged between variants, and the residual connection makes it a
near-identity transform at initialisation. Optional relay nodes representing transit
switches carry zero features and act purely as message-passing intermediaries,
propagating observations two hops through the topology.

The consequence for the attacks is structural rather than one of degree: without the
encoder an agent's policy reads its own observation alone, whereas with it a
perturbation applied to one agent reaches its neighbours' encoded states before
their actors. The single-agent differentiable path assumed by the critic-grounded
attack therefore does not hold for these variants, which must be attacked with the
joint formulation of Section~\ref{subsec:coordinated_attack} or excluded from that
comparison.

% ─────────────────────────────────────────────────────────────────────────────
\subsection{Variant matrix}
\label{subsec:variant_matrix}

\begin{table}[htbp]
  \centering
  \caption{MADDPG victim variants under test, named
           \texttt{<critic domain>-<value head>[-GNN]}.}
  \label{tab:variant_matrix}
  \begin{tabular}{llll}
    \hline
    \textbf{Variant} & \textbf{Critic} & \textbf{Value head} & \textbf{GNN} \\
    \hline
    CC-Simple        & Central & Simple   & --- \\
    CC-Duelling      & Central & Duelling & --- \\
    CC-Simple-GNN    & Central & Simple   & Yes \\
    CC-Duelling-GNN  & Central & Duelling & Yes \\
    LC-Simple        & Local   & Simple   & --- \\
    LC-Duelling      & Local   & Duelling & --- \\
    LC-Duelling-GNN  & Local   & Duelling & Yes \\
    \hline
  \end{tabular}
\end{table}

The matrix is not the full $2\times2\times2$ factorial: \texttt{LC-Simple-GNN} is
omitted, as the local-critic value-head comparison is already covered by
\texttt{LC-Simple}/\texttt{LC-Duelling} and the GNN comparison by
\texttt{LC-Duelling}/\texttt{LC-Duelling-GNN}. Where a single variant is reported,
\texttt{CC-Simple} serves as the reference configuration, being the canonical MADDPG
formulation with no additional architectural machinery.
